% 中药方剂模板
% 方剂模板.tex

\documentclass[12pt,UTF8]{ctexbook}

% 设置纸张信息。
\input{../../../Book/common/page}

% 设置字体，并解决显示难检字问题。
\xeCJKsetup{AutoFallBack=true}
% 注意：ctexbook类已默认设置SimSun为CJKrmdefault
% 以下设置用于确保字体回退和扩展字体可用
% 仅设置扩展字体以避免与默认设置冲突
\setCJKfamilyfont{hei}{SimHei}
\setCJKfamilyfont{kai}{KaiTi}

% 目录 chapter 级别加点（.）。
\usepackage{titletoc}
\titlecontents{chapter}[0pt]{\vspace{3mm}\bf\addvspace{2pt}\filright}{\contentspush{\thecontentslabel\hspace{0.8em}}}{}{\titlerule*[8pt]{.}\contentspage}

% 支持目录点击跳转
\usepackage[colorlinks,linkcolor=blue,citecolor=blue,urlcolor=blue]{hyperref}

% 设置 part 和 chapter 标题格式。
\ctexset{
	part/name= {第,卷},
	part/number={\chinese{part}},
	chapter/name={第,篇},
	chapter/number={\arabic{chapter}}
}

% 图片相关设置。
\usepackage{graphicx}
\graphicspath{{Images/}}

% 设置署名格式。
\newenvironment{shuming}{\hfill\zihao{4}}

% 注脚每页重新编号，避免编号过大。
\usepackage[perpage]{footmisc}

% 设置古文原文格式。
\newenvironment{yuanwen}{\bfseries\zihao{4}}

% 设置段落间距
\setlength{\parskip}{0.8em plus 0.2em minus 0.1em}

% 列表格式支持，列表项向右偏移。
\usepackage{enumitem}

% 定义古籍风格药物命令（行内下标格式，如：白术土炒₃两(90g)）
\newcommand{\gujiyaowusub}[2]{\textbf{#1}\raisebox{0.3ex}{\textsubscript{\scriptsize #2}}}

% 定义古籍风格药物表格环境（自动分行，上下对齐）
\newenvironment{gujiyaowutable}[1]{%
    \section{#1}%
    \begin{center}%
    \begin{tabular}{@{}p{3.5cm}p{3.5cm}p{3.5cm}p{3.5cm}@{}}%
}{%  % 4列布局，每列3.5cm宽度
    \end{tabular}%
    \end{center}%
}

\title{\heiti\zihao{0} 方剂名称}
\author{WangFei}
\date{\today}

\begin{document}

\maketitle
\tableofcontents

\frontmatter
\chapter{前言}

\mainmatter

\chapter{方剂名称}

\section{基本信息}
\begin{itemize}
	\item 【别名】：其他名称
	\item 【出处】：古籍或现代标准
	\item 【分类】：功效分类
\end{itemize}

\section{组成}

% 药物组成（古籍表格风格）
\begin{gujiyaowutable}{组成}
    \gujiyaowusub{白术土炒}{三两(90g)} & \gujiyaowusub{白芍炒}{二两(60g)} & \gujiyaowusub{陈皮炒}{一两半(45g)} & \gujiyaowusub{防风}{二两(60g)} \\
    \gujiyaowusub{人参}{一两(30g)} & \gujiyaowusub{茯苓}{一两(30g)} & \gujiyaowusub{甘草炙}{五钱(15g)} & \\
\end{gujiyaowutable}

\section{用法}
\begin{itemize}
	\item 【煎服方法】：煎煮时间、用水量
	\item 【服用剂量】：成人用量、儿童用量
	\item 【服用时间】：饭前/饭后、频服/顿服
	\item 【特殊处理】：先煎、后下、包煎等
\end{itemize}

\section{功效}
\begin{itemize}
	\item 【主要功效】：简明扼要的功效描述
	\item 【作用机制】：现代医学解释（可选）
\end{itemize}

\section{主治}
\begin{itemize}
	\item 【传统主治】：古籍记载的适应症
	\item 【现代主治】：现代临床应用范围
\end{itemize}

\section{方解}
\begin{description}
	\item[君药]：药物名称 - 作用
	\item[臣药]：药物名称 - 作用
	\item[佐药]：药物名称 - 作用
	\item[使药]：药物名称 - 作用
\end{description}

\section{临床应用}
\begin{itemize}
	\item 【辨证要点】：使用方剂的辨证依据
	\item 【现代应用】：现代医学疾病应用
	\item 【病例举隅】：典型病例（可选）
\end{itemize}

\section{加减变化}
\subsection{随证加减}
\begin{enumerate}[leftmargin=2cm]
	\item 
	\item 
	\item 
\end{enumerate}

\section{注意事项}
\subsection{使用禁忌}
\begin{enumerate}[leftmargin=2cm]
	\item 禁忌人群和情况
	\item 服用期间的注意事项
	\item 可能的副作用
\end{enumerate}

\section{附方}
\begin{itemize}
	\item 【相关方剂】：衍生方或类似方
	\item 【比较分析】：与相关方剂的区别
\end{itemize}

\backmatter

\end{document}